\chapter{マーケットメイクと高頻度取引}
\label{ch:mm}

ここまでの章は、数日〜数か月の予測を前提にしていた。この章は、秒以下の世界、つまり取引所の板の上で何が起きているかを扱う。ここで使う道具（板、スプレッド、逆選択、在庫）は、長い期間の運用の取引コスト（\chapref{ch:cost}）を理解するうえでも基本になる。

\section{板の仕組み}

\subsection{指値の注文と成行の注文}

多くの取引所は、\textbf{指値の注文の板}（リミット・オーダー・ブック）で売買を成立させる。

\begin{itemize}
\item \textbf{指値の注文}：「この価格以下でなら 100 株買う」。すぐには成立せず、板に並んで待つ。流動性を\textbf{供給}する。
\item \textbf{成行の注文}：「いますぐ 100 株買う」。板に並んでいる最も安い売りの指値と成立する。流動性を\textbf{消費}する。
\item \textbf{最良気配}：最も高い買いの指値（ビッド）と、最も安い売りの指値（アスク）。その差が\textbf{スプレッド}、中間が\textbf{仲値}。
\item \textbf{優先の順位}：同じ価格の指値の中では、先に出した注文から成立する（価格優先・時間優先）。良い順番を取るために速さが要る。
\item \textbf{ティック}：価格の最小の刻み。米国株では原則 1 セント。
\end{itemize}

\begin{figure}[H]
\centering
\begin{tikzpicture}[x=7mm,y=5mm]
\foreach \p/\q/\v in {1/3/99.95,2/5/99.96,3/2/99.97,4/6/99.98,5/4/99.99} {
  \fill[c3!55] ({-\q},{\p-0.4}) rectangle (0,{\p+0.4});
  \node[lab,anchor=west] at (0.15,\p) {\v};
}
\foreach \p/\q in {7/2,8/4,9/3,10/6,11/5} {
  \fill[c2!55] (0,{\p-0.4}) rectangle (\q,{\p+0.4});
  \node[lab,anchor=east] at (-0.15,\p) {100.0\the\numexpr \p-6\relax};
}
\draw[ink2,dashed] (-7,6) -- (7,6);
\node[lab,anchor=west] at (2.4,6.35) {スプレッド（100.01 − 99.99 = 0.02）};
\node[font=\footnotesize,c3!70!black] at (-4.5,0) {買いの指値（ビッド）};
\node[font=\footnotesize,c2!80!black] at (4.5,12.2) {売りの指値（アスク）};
\draw[arr] (0,-0.8)--(0,12.8) node[above,font=\scriptsize]{価格};
\end{tikzpicture}
\caption{指値の注文の板の模式図。棒の長さは、その価格に並んでいる注文の量。成行の買いは、最良の売り（100.01）から順に成立する。}
\label{fig:lob}
\end{figure}

\subsection{取引所の手数料の仕組み}

米国の多くの取引所は、流動性を供給した側（メイカー）に報奨金を払い、消費した側（テイカー）から手数料をとる（メイカー・テイカー方式）。マーケットメイカーの収益には、この報奨金も含まれる。

\section{マーケットメイクの経済学}

\subsection{収益と損失}

マーケットメイカーは、ビッドとアスクの両方に指値を出し続ける。買いと売りが同じ量ずつ来れば、1 往復ごとにスプレッドを稼げる。しかし、損失の源が二つある。

\begin{enumerate}
\item \textbf{逆選択}：相手が、自分の知らない情報を持っていることがある。良い知らせを知る相手は買い、マーケットメイカーは売ったあとで価格が上がる。
\item \textbf{在庫のリスク}：買いと売りが偏ると、在庫（ポジション）がたまり、その間の価格の変動のリスクを負う。
\end{enumerate}

\subsection{逆選択のモデル}

Glosten・Milgrom（1985）\cite{glostenmilgrom1985} のモデルでは、注文の一部は情報を持った投資家から、残りは情報を持たない投資家（流動性の必要から売買する）から来る。マーケットメイカーは、相手が誰かは分からないが、「買い注文が来た」こと自体が、価格が上がる可能性を示すと考えて、アスクを高めに、ビッドを低めに置く。スプレッドは、情報を持った投資家に負ける分を、持たない投資家から取り返す大きさに決まる。情報を持った投資家が多い銘柄（小型株、決算の前後）ほど、スプレッドは広い。

Kyle（1985）\cite{kyle1985} のモデルでは、情報を持った投資家は、自分の情報が価格に漏れすぎないように、少しずつ売買する。価格は、注文の流れ（買い − 売り）に比例して動く（$\Delta p = \lambda \times$ 注文の流れ）。$\lambda$ が大きい銘柄ほど、売買の価格への影響（インパクト）が大きい。

\subsection{在庫のモデル}

Ho・Stoll（1981）\cite{hostoll1981} と Avellaneda・Stoikov（2008）\cite{avellanedastoikov2008} は、在庫のリスクを考えて、指値の位置を決めるモデルを作った。Avellaneda・Stoikov のモデルでは、仲値を $s$、在庫を $q$（買い越しなら正）、リスクを嫌う度合いを $\gamma$、価格の変動を $\sigma$、残り時間を $T-t$ として、マーケットメイカーにとっての「本当の価格」（留保価格）は
\begin{equation}
r = s - q\,\gamma\,\sigma^2\,(T - t)
\end{equation}
である。在庫を多く持つほど、留保価格を下げ、ビッドもアスクも下に動かして、売りが成立しやすく、買いが成立しにくくする。こうして、在庫をゼロに戻す。スプレッドの幅は、リスクを嫌う度合いと、注文の来る頻度で決まる。

\section{高頻度取引}

\subsection{種類}

高頻度取引（HFT）は、計算機で秒以下の速さで注文を出し、取り消し、多くの売買をする取引の総称である。主な戦略は次のとおり。

\begin{itemize}
\item \textbf{電子的なマーケットメイク}：上で述べた、流動性の供給。HFT の最大の部分。
\item \textbf{遅延の裁定}：同じ資産（例：株価指数の先物と、その現物の ETF）が別の市場で取引されているとき、一方の価格の動きをもう一方より速く知って売買する。
\item \textbf{短期の統計的裁定}：板の偏り（買いの注文が売りより多い）などから、数秒〜数分の価格の動きを予測する。Cont・Kukanov・Stoikov（2014）\cite{cks2014} は、最良気配の注文の増減の偏り（注文の流れの不均衡）が、短い時間の価格の変化をよく説明することを示した。
\item \textbf{イベントの裁定}：経済指標や企業の発表を、機械で読んで瞬時に売買する。
\end{itemize}

\subsection{速さの競争}

取引所の計算機と同じ建物に自社の計算機を置く（コロケーション）、都市の間をマイクロ波の通信で結ぶ、注文の処理を専用の回路（FPGA）でする、などの投資で、何マイクロ秒を争う。Budish・Cramton・Shim（2015）\cite{bcs2015} は、連続的に売買を成立させる仕組みそのものが、この「速さの軍拡競争」を生むと論じ、一定の間隔（例：0.1 秒）ごとにまとめて売買を成立させる方式（頻繁なバッチ・オークション）を提案した。

\subsection{HFT の功罪}

\begin{itemize}
\item \textbf{良い面}：スプレッドを狭め、流動性を増やし、価格の発見を速くした、という実証が多い。
\item \textbf{悪い面}：2010 年 5 月 6 日の米国株の急落（フラッシュ・クラッシュ）のように、市場が混乱すると、HFT は指値を取り消して退き、流動性が一瞬で消えることがある。遅延の裁定は、遅い参加者から速い参加者へ利益を移すだけで、社会的な価値は小さいとも論じられる。
\end{itemize}

\section{個人や研究で扱えるか}

マーケットメイクと HFT は、ヘッジファンドの戦略の中で、個人や大学の研究が実際に競うことの最も難しい分野である。

\begin{itemize}
\item 実際の板のデータ（全注文の記録）は高価で、量が膨大である。
\item 競争の相手は、速さに巨額の投資をした専門の会社で、遅い参加者は逆選択の相手になる。
\item バックテストでは、自分の注文が板の何番目に並び、いつ成立するかを正確に再現できない。
\end{itemize}

一方で、この分野の考え方（スプレッド、逆選択、インパクト）は、遅い戦略の取引コストを見積もるうえで欠かせない。決算発表の直後の数分は、情報を持った投資家と HFT が最も活発に動く時間帯であり、LLM の読み取りで「当日に」売買しようとすると、この人たちと競うことになる。

\begin{projbox}
LLM-BLM の実験では、決算発表文の情報は発表の当日に株価に入りきった。その「入りきる」過程を担っているのが、まさにこの章の参加者（発表の文章を機械で読んで瞬時に売買する会社と、マーケットメイカー）である。学生の研究が発表の数分後に読み取りを終えても、その時点で情報の大部分は価格に入っている。研究の戦略を、「速さでは競わない」領域（数週〜数か月の、注目の低い銘柄の反応の遅れ）に置く必要があるのは、このためである。
\end{projbox}
